% ==============================================================================
% MÓDULO: CAPÍTULO III - ABORDAJE TEÓRICO (ENFOQUE CUALITATIVO FENOMENOLÓGICO)
% Archivo: subcapitulos/sec_03_abordaje_teorico.tex
% ==============================================================================
\section{ABORDAJE TEÓRICO}
\label{sec:cap3_teoria}

\subsection{Antecedentes de estudio}
\label{sec:3.1_antecedentes}
En estricta observancia del estándar metodológico cualitativo y de la cátedra de EN 486 (UNSCH), a cargo del Dr. Manglio Aguirre Andrade, la revisión de la literatura se orienta a comprender cómo se ha investigado la \textbf{experiencia vivida y el significado de la calidad asistencial y del encuentro interpersonal}:

\subsubsection{Antecedentes internacionales}
\label{sec:3.1.1_internacionales}
\begin{itemize}
    \item \textbf{Gómez-Restrepo y colaboradores (2023, Colombia):} En su investigación cualitativa fenomenológica titulada \textit{``Experiencia del usuario sobre la calidad del trato interpersonal en el primer nivel de atención''}, analizaron 18 entrevistas en profundidad en centros de salud primarios. Concluyeron que la percepción del usuario está gobernada por la calidez comunicativa, el contacto visual y la empatía del profesional de enfermería. Demostraron que una relación de cuidado humanizada genera bienestar emocional y confianza terapéutica, independientemente de las limitaciones materiales del entorno asistencial. \textit{Aporte al estudio:} Aporta categorías directas sobre la validación afectiva del paciente en la consulta.
    \item \textbf{Pérez y Silva (2022, Ecuador):} En su estudio cualitativo hermenéutico \textit{``Significados de la relación enfermera-paciente desde la vivencia del usuario andino''}, entrevistaron a 15 usuarios quechuahablantes. Sus hallazgos revelaron que el saludo cordial en lengua originaria, el tono de voz pausado y la escucha libre de juicios son interpretados por los usuarios como los máximos indicadores de dignidad y respeto profesional. \textit{Aporte al estudio:} Provee sustento epistemológico para incorporar la pertinencia quechua chanka en el protocolo de contacto.
    \item \textbf{Martínez-Salazar y Morales (2021, México):} En su investigación cualitativa \textit{``Humanización del cuidado y relación interpersonal en la consulta ambulatoria''}, desarrollada con 20 usuarios en centros comunitarios, identificaron que el uso de tecnicismos incomprensibles y la prisa del profesional provocan sentimientos de intimidación y silencio en el paciente, impidiéndole manifestar sus verdaderos temores. \textit{Aporte al estudio:} Justifica la exploración de la pedagogía comunicativa del personal de salud.
\end{itemize}

\subsubsection{Antecedentes nacionales}
\label{sec:3.1.2_nacionales}
\begin{itemize}
    \item \textbf{Mamani y Quispe (2024, Lima):} En su investigación cualitativa \textit{``Percepción del trato humanizado en el cuidado de enfermería en el primer nivel de atención''}, determinaron mediante 16 relatos vivenciales que la frialdad en la mirada, la falta de saludo inicial y las distracciones del personal durante la consulta son valoradas por el usuario como una despersonalización del acto de cuidado. \textit{Aporte al estudio:} Establece pautas para auditar los distractores y la presencia atenta del profesional en el consultorio.
    \item \textbf{Córdova y Reyes (2023, Arequipa):} En su estudio cualitativo \textit{``Empatía, comunicación y percepción del cuidado en la atención primaria''}, concluyeron que la experiencia del usuario se transforma positivamente cuando la enfermera o médico dedican tiempo a indagar por el estado emocional del paciente y no únicamente por sus síntomas biológicos. \textit{Aporte al estudio:} Fundamenta la subcategoría de contención afectiva y empatía.
\end{itemize}

\subsubsection{Antecedentes regionales y locales (UNSCH -- Ayacucho)}
\label{sec:3.1.3_locales}
\begin{itemize}
    \item \textbf{Velarde Quispedina (2023, Ayacucho -- UNSCH):} En su investigación sobre la atención asistencial en Huamanga \cite{velarde2023}, evidenció que el encuentro interpersonal y el trato brindado por enfermería constituyen el núcleo vivencial determinante del bienestar del paciente, superando con creces la relevancia de las condiciones materiales y edilicias del entorno. \textit{Aporte al estudio:} Evidencia regional sobre la primacía ontológica del trato humano frente al espacio físico.
    \item \textbf{Ruiz (2024, Ayacucho -- UNSCH):} En su investigación en el Hospital de Apoyo de Huanta \cite{ruiz2024}, comprobó que la empatía y la cohesión interna del equipo asistencial se proyectan de manera directa en la calidez, la paciencia y el respeto brindado al usuario externo. \textit{Aporte al estudio:} Sustenta la articulación entre la actitud del profesional y la vivencia sentida del usuario.
    \item \textbf{Aguirre-Andrade y colaboradores (2020, Ayacucho -- UNSCH):} En su estudio científico en comunidades de Cangallo \cite{aguirre2020}, demostraron que la interacción asistencial en lengua originaria (quechua) y el respeto irrestricto a la cosmovisión andina son factores insustituibles para el éxito del cuidado comunitario. \textit{Aporte al estudio:} Consolida el pilar de adecuación intercultural en la interacción del cuidado.
\end{itemize}

\subsection{Bases teóricas fundamentales}
\label{sec:3.2_base_teorica}

\subsubsection{Dimensión del Proceso Interpersonal en la Teoría de Avedis Donabedian}
\label{sec:3.2.1_donabedian}
En el modelo conceptual de Avedis Donabedian \cite{donabedian1980, donabedian2005}, la calidad asistencial comprende tres esferas: Estructura, Proceso y Resultado. El presente estudio adopta como fundamento cualitativo central el \textbf{Componente de Proceso en su vertiente interpersonal}:
\begin{quote}
``El manejo de la relación interpersonal es el vehículo a través del cual la atención técnica se aplica y se hace efectiva. Comprende la calidez, la empatía, el respeto a la dignidad del paciente, la privacidad, la comunicación clara y la ética demostrada en el contacto directo'' \cite{donabedian2005}.
\end{quote}
Donabedian enfatiza que la pericia técnica carece de impacto terapéutico pleno si el acto de interacción fractura la confianza y la dignidad del usuario.

\subsubsection{Teoría del Cuidado Humano y Factores Caritas de Jean Watson}
\label{sec:3.2.2_watson}
La teorista Jean Watson \cite{watson2008} define el cuidado como el núcleo ontológico, existencial y ético de la enfermería. Postula que una atención de auténtica calidad constituye un \textbf{encuentro intersubjetivo transpersonal} guiado por los Factores Caritas:
\begin{enumerate}
    \item Cultivo de la amabilidad, la sensibilidad y el respeto irrestricto a la dignidad humana.
    \item Construcción de una relación auténtica de ayuda-confianza mediante la escucha activa y empática.
    \item Aceptación incondicional de la expresión de sentimientos positivos y negativos del paciente doliente.
    \item Creación de un entorno de cuidado protector que resguarde el pudor físico, mental y espiritual.
\end{enumerate}

\subsubsection{Fundamentos del Cuidado Fenomenológico en Enfermería: Patricia Benner y Holly Wilson}
\label{sec:3.2.3_benner_wilson}
En concordancia con los estudios fenomenológicos empíricos en enfermería destacados por Norlyk y Harder \cite{norlyk2010} y citados por Hernández-Sampieri y Mendoza \cite{sampieri2018}:
\begin{itemize}
    \item \textbf{Patricia Benner (2008) \cite{benner2008}:} Sostiene que el cuidado en enfermería es una práctica relacional, ética e intuitiva que no puede mecanizarse. En la fenomenología clínica, el profesional y el investigador comprenden que el paciente no es un objeto con síntomas aislados, sino un ser humano inmerso en una situación vivencial de vulnerabilidad.
    \item \textbf{Holly Wilson (2007) \cite{wilson2007}:} Señala que la fenomenología empírica en salud tiene como misión develar la esencia de las vivencias del paciente, suspendiendo los presupuestos teóricos del observador para entender qué significa experimentar la atención sanitaria desde la perspectiva de quien la recibe.
\end{itemize}

\subsection{Matriz de categorización cualitativa y definición de unidades de significado}
\label{sec:3.3_categorias}
De acuerdo con Roberto Hernández-Sampieri, Carlos Fernández-Collado y Pilar Baptista-Lucio (2014, Cap. 14 y 15) \cite{sampieri2018}, en la investigación fenomenológica se sustituye la tradicional operacionalización positivista de variables por una \textbf{Matriz de Categorización Cualitativa}, la cual articula los conceptos potenciales de partida con las cuatro dimensiones existenciales del mundo de la vida (\textit{Lebenswelt}) y las categorías emergentes (esenciales y divergentes):

\begin{table}[H]
\centering
\scriptsize
\begin{tabularx}{\textwidth}{p{2.3cm}p{2.5cm}p{3.2cm}p{3.8cm}X}
\toprule
\textbf{Concepto Potencial} & \textbf{Dimensión Existencial} & \textbf{Definición Conceptual / Significado Vivencial} & \textbf{Categorías Emergentes y Códigos Inductivos} & \textbf{Pregunta Guía de Profundización} \\
\midrule
\multirow{5}{=}{\textbf{CALIDAD DE ATENCIÓN Y CUIDADO HUMANO}\newline\newline \textit{(Encuentro intersubjetivo y vivencia del usuario en el C.S. Belén)}} 
& \textbf{1. Temporalidad (Tiempo vivido)} 
& Duración efectiva del contacto clínico en consultorio, ritmo pausado vs. prisa protocolar y dedicación de tiempo para escuchar y explicar con calma. 
& $\bullet$ \textit{Esencia compartida:} Desazón ante consultas apresuradas de escasos minutos; bienestar profundo cuando el profesional no muestra apuro y dialoga con paciencia.\newline $\bullet$ \textit{Variaciones individuales:} Madres en CRED priorizan el tiempo dedicado a la consejería; adultos mayores demandan calma en la explicación de sus recetas. 
& ¿Cómo experimentó el tiempo que el personal de salud le dedicó en la consulta? ¿Sintió que la atención fue apresurada o con la calma necesaria? \\
\cmidrule{2-5}
& \textbf{2. Espacialidad (Espacio vivido)} 
& Vivencia de la privacidad acústica y visual dentro del consultorio, confort físico de las salas y resguardo estricto de la confidencialidad clínica. 
& $\bullet$ \textit{Esencia compartida:} Incomodidad e inhibición ante tabiquerías delgadas que filtran voces; tranquilidad y confianza en ambientes cerrados y dignos.\newline $\bullet$ \textit{Variaciones individuales:} Consultorios de Obstetricia y CRED demandan mayor privacidad física que salas de triaje general. 
& ¿Sintió que el consultorio le brindó la intimidad, reserva y confidencialidad necesaria para dialogar y ser examinado(a)? \\
\cmidrule{2-5}
& \textbf{3. Corporalidad (Cuerpo vivido)} 
& Sentires somáticos durante el acto de cuidado: respeto al pudor físico al desvestirse, delicadeza en procedimientos y contención del dolor somático (\textit{nanay}). 
& $\bullet$ \textit{Esencia compartida:} Confianza y alivio ante el tacto suave y delicado en curaciones e inyecciones; consuelo humano frente al cuerpo adolorido (\textit{onqoy}).\newline $\bullet$ \textit{Variaciones individuales:} Madres sensibles al trato corporal hacia sus hijos; adultos mayores valoran la gentileza física ante sus limitaciones motrices. 
& Al ser examinado(a) físicamente o recibir curaciones, ¿cómo sintió que el personal trató su cuerpo y su dolor somático? \\
\cmidrule{2-5}
& \textbf{4. Relacionalidad (Relación humana vivida)} 
& Encuentro intersubjetivo entre el usuario y el equipo asistencial: calidez de la mirada, acogida cordial (\textit{allin chaskiy}), respeto mutuo (\textit{respetanakuy}) y lengua quechua. 
& $\bullet$ \textit{Esencia compartida:} Primacía del contacto visual sobre la pantalla del computador; consuelo al ser escuchado con ternura y en quechua chanka.\newline $\bullet$ \textit{Variaciones individuales:} Reconocimiento de alta calidez en Enfermería frente a mayor prisa en turnos médicos de alta demanda. 
& ¿Cómo fue el trato interpersonal y el encuentro humano con el equipo de salud? ¿Sintió calidez, mirada atenta y escucha en su lengua originaria? \\
\cmidrule{2-5}
& \textbf{5. Esencia Holística de la Calidad} 
& Significado integral y valorativo que el usuario otorga al cuidado humanizado, digno y resolutivo en el primer nivel asistencial. 
& $\bullet$ \textit{Esencia compartida:} Calidad es ser reconocido como persona humana integral, conjugando destreza técnica con calidez afectiva incondicional.\newline $\bullet$ \textit{Variaciones individuales:} Adultos mayores priorizan comprensión y paciencia; personas jóvenes valoran la claridad diagnóstica y la empatía en el trato. 
& Al reflexionar sobre la atención recibida, ¿qué significa para usted recibir un cuidado de salud con auténtica calidad humana? \\
\bottomrule
\end{tabularx}
\caption{Matriz de categorización cualitativa y definición de unidades de significado según Hernández-Sampieri.}
\label{tab:categorizacion}
\end{table}

\newpage
